\begin{tomtat}
	\begin{minipage}[t]{0.3\textwidth}
		\vspace{0pt}
		\begin{tomtat}
			\begin{center}
				\textbf{ĐỀ 31}
			\end{center}
			
		\end{tomtat}
	\end{minipage}
	\hfill
	\begin{minipage}[t]{0.5\textwidth}
		\vspace{0pt}
		\begin{tomtat}
			\begin{center}
				\includegraphics[scale=0.08]{images/31}
			\end{center}
		\end{tomtat}
	\end{minipage}
\end{tomtat}
\textbf{I. PHẦN TRẮC NGHIỆM}
\begin{ex}
	Giá trị của biểu thức $A=\sqrt[3]{-8}+\sqrt{36}$ là:
	\choice{$-8$}{$8$}{$-4$}{$4$}
	\loigiai{
		Ta có:
		\begin{align*} A &= \sqrt[3]{-8} + \sqrt{36} \\ &= \sqrt[3]{(-2)^3} + \sqrt{6^2} \\ &= -2 + 6 \\ &= 4 \end{align*}
		\\ \textbf{Đáp Án D}
	}
\end{ex}

\begin{ex}
	Nghiệm của phương trình: $3x – 5 = 7$ là:
	\choice{$-4$}{$4$}{$-5$}{$5$}
	\loigiai{
		Ta có: $3x - 5 = 7$
		\begin{align*} 3x &= 7+5 \\ 3x &= 12 \\ x &= \frac{12}{3} \\ x &= 4 \end{align*}
		\\ \textbf{Đáp Án B}
	}
\end{ex}

\begin{ex}
	Cho hàm số $y=f(x)=2x^2$. $f(-1)$ có giá trị là:
	\choice{$1$}{$-1$}{$2$}{$-2$}
	\loigiai{
		Ta có $f(x) = 2x^2$.
		Khi $x=-1$, thì $f(-1) = 2(-1)^2 = 2 \cdot 1 = 2$.
		\\ \textbf{Đáp Án C}
	}
\end{ex}

\begin{ex}
	Gieo một con xúc xắc 60 lần cho kết quả như sau:
	\begin{center}
		\begin{tabular}{|c|c|c|c|c|c|c|}
			\hline
			Số chấm& 1 & 2 & 3 & 4 & 5 & 6 \\
			\hline
			Tần số & 10 & ? & 9 & 8 & 12 & 11 \\
			\hline
		\end{tabular}
	\end{center}
	Tần số xuất hiện mặt 2 chấm là:
	\choice{$7$}{$8$}{$9$}{$10$}
	\loigiai{
		Tổng số lần gieo là 60.
		Tổng tần số của các mặt đã biết là: $10+9+8+12+11 = 50$.
		Gọi tần số xuất hiện mặt 2 chấm là $n_2$.
		Ta có: $10 + n_2 + 9 + 8 + 12 + 11 = 60$
		$n_2 + 50 = 60$
		$n_2 = 60 - 50$
		$n_2 = 10$.
		\\ \textbf{Đáp Án D}
	}
\end{ex}

\begin{ex}
	Nghiệm của bất phương trình: $5x-10 \ge 0$ là:
	\choice{$x \ge 2$}{$x \ge -2$}{$x \le 2$}{$x \le -2$}
	\loigiai{
		Ta có: $5x-10 \ge 0$
		\begin{align*} 5x &\ge 10 \\ x &\ge \frac{10}{5} \\ x &\ge 2 \end{align*}
		\\ \textbf{Đáp Án A}
	}
\end{ex}
\begin{ex}
	Điều kiện của $x$ để biểu thức $A = \sqrt{x-2023}$ có nghĩa là:
	\choice{$x \ge 2023$}{$x > 2023$}{$x < 2023$}{$x \le 2023$}
	\loigiai{
		Biểu thức $A = \sqrt{x-2023}$ có nghĩa khi $x-2023 \ge 0$.
		$x-2023 \ge 0 \Leftrightarrow x \ge 2023$.
		\\ \textbf{Đáp Án A}
	}
\end{ex}

\begin{ex}
	Hàm số nào sau đây là hàm số bậc nhất?
	\choice{$y = 2x+1$}{$y = x^2+1$}{$y = \frac{1}{x}$}{$y = \sqrt{x}$}
	\loigiai{
		Hàm số bậc nhất có dạng $y = ax+b$ với $a \neq 0$.
		A. $y = 2x+1$ là hàm số bậc nhất ($a=2, b=1$).
		B. $y = x^2+1$ là hàm số bậc hai.
		C. $y = \frac{1}{x}$ không phải là hàm số bậc nhất (hàm phân thức).
		D. $y = \sqrt{x}$ không phải là hàm số bậc nhất (hàm căn thức).
		\\ \textbf{Đáp Án A}
	}
\end{ex}

\begin{ex}
	Trong các hệ phương trình sau, hệ phương trình nào là hệ hai phương trình bậc nhất hai ẩn?
	\choice{$\begin{cases} x+y=1 \\ x-y=2 \end{cases}$}{$\begin{cases} x^2+y=1 \\ x-y=2 \end{cases}$}{$\begin{cases} x+y^2=1 \\ x-y=2 \end{cases}$}{$\begin{cases} x+y=1 \\ xy=2 \end{cases}$}
	\loigiai{
		Hệ hai phương trình bậc nhất hai ẩn có dạng $\begin{cases} a_1x+b_1y=c_1 \\ a_2x+b_2y=c_2 \end{cases}$, trong đó $a_1, b_1$ không đồng thời bằng 0 và $a_2, b_2$ không đồng thời bằng 0.
		A. $\begin{cases} x+y=1 \\ x-y=2 \end{cases}$ là hệ hai phương trình bậc nhất hai ẩn.
		B. $\begin{cases} x^2+y=1 \\ x-y=2 \end{cases}$ có phương trình $x^2+y=1$ là bậc hai theo $x$.
		C. $\begin{cases} x+y^2=1 \\ x-y=2 \end{cases}$ có phương trình $x+y^2=1$ là bậc hai theo $y$.
		D. $\begin{cases} x+y=1 \\ xy=2 \end{cases}$ có phương trình $xy=2$ không phải là bậc nhất.
		\\ \textbf{Đáp Án A}
	}
\end{ex}

\begin{ex}
	Cho tam giác ABC vuông tại A, đường cao AH. Hệ thức nào sau đây sai?
	\choice{$AB^2 = BH \cdot BC$}{$AC^2 = CH \cdot BC$}{$AH^2 = BH \cdot CH$}{$AB \cdot AC = AH \cdot CH$}
	\loigiai{
		Trong tam giác vuông ABC, đường cao AH:
		A. $AB^2 = BH \cdot BC$ (Hệ thức lượng trong tam giác vuông) - Đúng.
		B. $AC^2 = CH \cdot BC$ (Hệ thức lượng trong tam giác vuông) - Đúng.
		C. $AH^2 = BH \cdot CH$ (Hệ thức lượng trong tam giác vuông) - Đúng.
		D. $AB \cdot AC = AH \cdot BC$ (Hệ thức $bc=ah$) chứ không phải $AB \cdot AC = AH \cdot CH$. Vậy D sai.
		\\ \textbf{Đáp Án D}
	}
\end{ex}

\begin{ex}
	Với giá trị nào của $m$ thì hàm số $y=(m-2)x+3$ đồng biến trên $\mathbb{R}$?
	\choice{$m=2$}{$m>2$}{$m<2$}{$m \neq 2$}
	\loigiai{
		Hàm số $y=(m-2)x+3$ là hàm số bậc nhất khi $m-2 \neq 0 \Leftrightarrow m \neq 2$.
		Hàm số bậc nhất $y=ax+b$ đồng biến trên $\mathbb{R}$ khi $a>0$.
		Do đó, hàm số $y=(m-2)x+3$ đồng biến trên $\mathbb{R}$ khi $m-2 > 0 \Leftrightarrow m > 2$.
		\\ \textbf{Đáp Án B}
	}
\end{ex}
\newpage
\textbf{II. PHẦN TỰ LUẬN}
\begin{ex}
	Giải hệ phương trình: \\
	$\begin{cases} 2x - y = 2 \\ 9x + 8y = 34 \end{cases} $
	\loigiai{
		Ta có hệ phương trình:
		$$ \begin{cases} 2x - y = 2 \quad &(1) \\ 9x + 8y = 34 \quad &(2) \end{cases} $$
		Từ phương trình (1), suy ra $y = 2x - 2$.
		Thế $y = 2x - 2$ vào phương trình (2):
		\begin{align*} 9x + 8(2x - 2) &= 34 \\ 9x + 16x - 16 &= 34 \\ 25x &= 34 + 16 \\ 25x &= 50 \\ x &= \frac{50}{25} \\ x &= 2 \end{align*}
		Thay $x=2$ vào $y = 2x - 2$:
		\begin{align*} y &= 2(2) - 2 \\ y &= 4 - 2 \\ y &= 2 \end{align*}
		Vậy nghiệm của hệ phương trình là $(x; y) = (2; 2)$.
		\\ \textbf{Đáp số: $(x;y) = (2;2)$}
	}
\end{ex}
\begin{ex}
	Rút gọn biểu thức: $\left( \frac{1}{2\sqrt{x}+1} + \frac{1}{2\sqrt{x}-1} \right) : \frac{1}{1-4x}$ với $x \ge 0; x \neq \frac{1}{4}$.
	\loigiai{
		Điều kiện: $x \ge 0; x \neq \frac{1}{4}$.
		Ta có:
		\begin{align*} A &= \left( \frac{1}{2\sqrt{x}+1} + \frac{1}{2\sqrt{x}-1} \right) : \frac{1}{1-4x} \\ &= \left( \frac{(2\sqrt{x}-1) + (2\sqrt{x}+1)}{(2\sqrt{x}+1)(2\sqrt{x}-1)} \right) \cdot (1-4x) \\ &= \frac{2\sqrt{x}-1 + 2\sqrt{x}+1}{(2\sqrt{x})^2 - 1^2} \cdot (1-4x) \\ &= \frac{4\sqrt{x}}{4x-1} \cdot (1-4x) \end{align*}
		Vì $1-4x = -(4x-1)$, và $x \neq \frac{1}{4}$ nên $4x-1 \neq 0$.
		\begin{align*} A &= \frac{4\sqrt{x}}{4x-1} \cdot (-(4x-1)) \\ &= \frac{4\sqrt{x}}{\cancel{(4x-1)}} \cdot (-\cancel{(4x-1)}) \\ &= -4\sqrt{x} \end{align*}
		Vậy $A = -4\sqrt{x}$.
		\\ \textbf{Đáp số: $A = -4\sqrt{x}$}
	}
\end{ex}
\begin{ex}
	Tìm m để phương trình $x^2 - 2x - 3 = 0$. Không giải phương trình, hãy tính giá trị biểu thức sau $T = x_1x_2 + x_1^2 + x_2^2$.
	\loigiai{
		Đề bài có vẻ như bị nhầm lẫn giữa yêu cầu "Tìm m" và "Không giải phương trình, hãy tính giá trị biểu thức". Giả sử đề bài yêu cầu tính giá trị biểu thức $T$ với phương trình $x^2 - 2x - 3 = 0$ đã cho.
		
		Phương trình $x^2 - 2x - 3 = 0$ có $a=1, b=-2, c=-3$.
		$\Delta' = b'^2 - ac = (-1)^2 - 1(-3) = 1+3=4 > 0$.
		Phương trình có hai nghiệm phân biệt $x_1, x_2$.
		Theo hệ thức Vi-ét, ta có:
		$$ \begin{cases} x_1 + x_2 = -\frac{b}{a} = -\frac{-2}{1} = 2 \\ x_1 x_2 = \frac{c}{a} = \frac{-3}{1} = -3 \end{cases} $$
		Ta biến đổi biểu thức T:
		\begin{align*} T &= x_1x_2 + x_1^2 + x_2^2 \\ &= x_1x_2 + (x_1+x_2)^2 - 2x_1x_2 \\ &= (x_1+x_2)^2 - x_1x_2 \end{align*}
		Thay giá trị từ hệ thức Vi-ét vào biểu thức T:
		\begin{align*} T &= (2)^2 - (-3) \\ &= 4 + 3 \\ &= 7 \end{align*}
		Vậy $T=7$.
		\\ \textbf{Đáp số: $T=7$}
	}
\end{ex}

\begin{ex}
	Trong kỳ thi Tuyển sinh vào lớp 10 THPT năm 2019, tổng chỉ tiêu tuyển sinh của Trường THPT A và trường THPT B là 900 học sinh. Do cả hai trường đều có chất lượng giáo dục rất tốt nên sau khi hết hạn thời gian điều chỉnh nguyện vọng thì số lượng thí sinh đăng ký dự tuyển vào Trường THPT A và Trường THPT B tăng lần lượt là 15\% và 10\% so với chỉ tiêu ban đầu. Vì vậy, tổng số thí sinh đăng ký dự tuyển của cả hai trường là 1010. Hỏi số lượng thí sinh đăng ký dự tuyển của mỗi trường là bao nhiêu?
	\loigiai{
		Gọi chỉ tiêu tuyển sinh ban đầu của Trường THPT A là $x$ (học sinh).
		Gọi chỉ tiêu tuyển sinh ban đầu của Trường THPT B là $y$ (học sinh).
		Điều kiện: $x, y \in \mathbb{N}^*$.
		
		Tổng chỉ tiêu tuyển sinh ban đầu của hai trường là 900 học sinh:
		$x + y = 900 \quad (1)$
		
		Số lượng thí sinh đăng ký dự tuyển vào Trường THPT A sau khi tăng 15\% là:
		$x + 0,15x = 1,15x$ (học sinh).
		Số lượng thí sinh đăng ký dự tuyển vào Trường THPT B sau khi tăng 10\% là:
		$y + 0,10y = 1,10y$ (học sinh).
		
		Tổng số thí sinh đăng ký dự tuyển của cả hai trường sau khi tăng là 1010:
		$1,15x + 1,10y = 1010 \quad (2)$
		
		Từ (1) và (2), ta có hệ phương trình:
		$$ \begin{cases} x + y = 900 \\ 1,15x + 1,10y = 1010 \end{cases} $$
		Giải hệ phương trình:
		Từ (1) $\Rightarrow y = 900 - x$.
		Thế $y = 900 - x$ vào (2):
		\begin{align*} 1,15x + 1,10(900 - x) &= 1010 \\ 1,15x + 990 - 1,10x &= 1010 \\ 0,05x &= 1010 - 990 \\ 0,05x &= 20 \\ x &= \frac{20}{0,05} \\ x &= 400 \end{align*}
		Thay $x = 400$ vào $y = 900 - x$:
		$y = 900 - 400 = 500$.
		(Thỏa mãn điều kiện $x, y \in \mathbb{N}^*$).
		
		Số lượng thí sinh đăng ký dự tuyển của Trường THPT A là: $1,15x = 1,15 \cdot 400 = 460$ học sinh.
		Số lượng thí sinh đăng ký dự tuyển của Trường THPT B là: $1,10y = 1,10 \cdot 500 = 550$ học sinh.
		Kiểm tra: $460 + 550 = 1010$ (đúng).
		
		Vậy, số lượng thí sinh đăng ký dự tuyển của Trường THPT A là 460 học sinh và của Trường THPT B là 550 học sinh.
	}
\end{ex}
\begin{ex}
	Một túi đựng 4 viên bi có cùng khối lượng và kích thước, được đánh số 1; 2; 3; 4. Lấy ngẫu nhiên lần lượt 2 viên bi từ túi đó, viên bi lấy ra lần đầu không trả lại vào túi. Mô tả không gian mẫu của phép thử và tính xác suất để lấy được 2 viên bi mà tổng hai số trên hai viên bi đó là số lẻ.
	\loigiai{
		Không gian mẫu $\Omega$:
		Lấy lần 1 được 1: $(1;2)$; $(1;3)$; $(1;4)$.
		Lấy lần 1 được 2: $(2;1)$; $(2;3)$; $(2;4)$.
		Lấy lần 1 được 3: $(3;1)$; $(3;2)$; $(3;4)$.
		Lấy lần 1 được 4: $(4;1)$; $(4;2)$; $(4;3)$.
		Vậy, $\Omega = \{ (1;2); (1;3); (1;4); (2;1); (2;3); (2;4); (3;1); (3;2); (3;4); (4;1); (4;2); (4;3) \}$.
		Số phần tử của không gian mẫu: $n(\Omega) = 12$.
		
		Gọi A là biến cố: “Lấy được 2 viên bi mà tổng hai số trên hai viên bi đó là số lẻ”.
		Để tổng hai số là số lẻ, một số phải chẵn và một số phải lẻ.
		Các cặp có tổng là số lẻ:
		\begin{itemize}
			\item Lần 1 lẻ, lần 2 chẵn: $(1;2)$; $(1;4)$; $(3;2)$; $(3;4)$.
			\item Lần 1 chẵn, lần 2 lẻ: $(2;1)$; $(2;3)$; $(4;1)$; $(4;3)$.
		\end{itemize}
		Vậy, $A = \{ (1;2); (1;4); (3;2); (3;4); (2;1); (2;3); (4;1); (4;3) \}$.
		Số phần tử của biến cố A: $n(A) = 8$.
		
		Xác suất của biến cố A:
		$P(A) = \frac{n(A)}{n(\Omega)} = \frac{8}{12} = \frac{2}{3}$.
	}
\end{ex}
\begin{ex}
	Cho tam giác ABC nhọn $AB < AC$ có đường cao AD và đường phân giác trong AO (D, O thuộc cạnh BC). Kẻ OM vuông góc với AB tại M, ON vuông góc với AC tại N.
	\begin{enumerate}
		\item[a)] Chứng minh bốn điểm A, D, O, N cùng nằm trên một đường tròn.
		\item[b)] Chứng minh: $\widehat{BDM} = \widehat{ODN}$. Qua O, kẻ đường thẳng vuông góc với BC cắt MN tại I, AI cắt BC tại K. Chứng minh K là trung điểm của BC.
	\end{enumerate}
\end{ex}


\setcounterref{ex}{0}
